% diagrams/firstwords/casa-vecinos.tex -- la casa de al lado, con los
% vecinos nuevos asomados a las ventanas: Amy y su madre arriba, Mr
% Brown y Pip abajo. Cada cabeza va recortada por su ventana.
\begin{tikzpicture}[dibujofw]
  % el tejado y la chimenea
  \filldraw[fill=white] (7.2,9.0) rectangle (8.2,10.6);
  \filldraw[fill=white] (-0.7,7.0) -- (1.6,9.8) -- (8.4,9.8) -- (10.7,7.0) -- cycle;
  % la fachada
  \filldraw[fill=white] (0,0) rectangle (10,7.0);
  % las ventanas, con quien se asoma a cada una
  \foreach \x/\y/\quien in {1.1/4.1/Amy, 6.4/4.1/MrsBrown, 1.1/0.9/MrBrown} {
    \filldraw[fill=white] (\x,\y) rectangle ++(2.5,2.3);
    \begin{scope}
      \clip (\x,\y) rectangle ++(2.5,2.3);
      \filldraw[fill=white, rounded corners=3mm] (\x+0.35,\y-0.3) rectangle ++(1.8,1.1);
      \begin{scope}[shift={(\x+1.25,\y+1.3)}, scale=0.62]
        \csname cabeza\quien\endcsname
      \end{scope}
    \end{scope}
    \draw (\x,\y) rectangle ++(2.5,2.3);
    \filldraw[fill=white] (\x-0.2,\y-0.25) rectangle ++(2.9,0.25);
  }
  % la ventana de Pip, con la jaula
  \filldraw[fill=white] (6.4,0.9) rectangle ++(2.5,2.3);
  \begin{scope}
    \clip (6.4,0.9) rectangle ++(2.5,2.3);
    \begin{scope}[shift={(7.65,1.0)}, scale=0.3]
      \jaulaFondoFW
      \begin{scope}[shift={(0,2.9)}, scale=0.62] \perchaFW{-3.2}{3.2} \pipFW \end{scope}
      \jaulaBaseFW
    \end{scope}
  \end{scope}
  \draw (6.4,0.9) rectangle ++(2.5,2.3);
  \filldraw[fill=white] (6.2,0.65) rectangle ++(2.9,0.25);
  % la puerta, con su pomo, y el escalón
  \filldraw[fill=white] (4.2,0.3) -- (4.2,3.2) arc[start angle=180, end angle=0, radius=0.8]
    -- (5.8,0.3) -- cycle;
  \filldraw[fill=white] (5.45,1.7) circle (0.12);
  \filldraw[fill=white] (3.9,0) rectangle (6.1,0.3);
  % el suelo
  \draw (-1.2,0) -- (11.2,0);
\end{tikzpicture}
